\Backup{regular Bohr sets and translation errors}{}{
\begin{block}{Regularity and regular radii}
Write $d=\max(1,|S|)$. A Bohr set $P=\Bohr(S,r)$ is regular if
\[(1-100d|\kappa|)|P|\leq|\Bohr(S,(1+\kappa)r)|
\leq(1+100d|\kappa|)|P|\]
for $|\kappa|\leq1/(100d)$. For every $0<\varepsilon<1$ there
exists $r\in[\varepsilon,2\varepsilon]$ with $\Bohr(S,r)$ regular.
\end{block}
\small If $P$ is regular, $|F|\leq1$, and $\norm{t}_S\leq\varepsilon r$
with $\varepsilon\leq1/(100d)$, then
\[\abs{\E_{x\in P}F(x+t)-\E_{x\in P}F(x)}\leq200d\varepsilon.\]
\source{\GT{Definition 2.6, Lemmas 8.2 and 4.2}. Existence of regular radii is a standard black box.}
}{
The case $S=\varnothing$ is trivial: the Bohr set is $G$. Otherwise the
definition agrees with the paper. The regular-radius existence theorem is
the auxiliary geometric input we use without proof.

The translation estimate follows directly: both $P$ and $P+t$ contain
$\Bohr(S,(1-\varepsilon)r)$ and are contained in
$\Bohr(S,(1+\varepsilon)r)$. Thus the symmetric difference has size
at most $200d\varepsilon|P|$. Dividing by $|P|$ proves the inequality.
For a small Bohr set $Q=\Bohr(S,s)$, the inclusion
$P+Q\subseteq\Bohr(S,r+s)$ also gives $|P+Q|\leq2|P|$ when
$s/r\leq1/(100d)$. These are precisely the boundary estimates used
in the symmetry and final Fourier steps.
}

\Backup{the finite-field model for Lemma 2.5}{}{
\small Take $G=\F_5^n$ and identify $\Gh$ with $\F_5^n$.
\begin{center}
\begin{tikzpicture}[node distance=6mm]
\node[box,text width=10cm] (a) {BSG: a large graph subset $\Gamma'$ has small doubling};
\node[box,text width=10cm,below=of a] (b) {Finite-field Freiman: $\Gamma'\subseteq W\leq\F_5^{2n}$,\quad $|W|\leq C_\eta|G|$};
\node[box,text width=10cm,below=of b] (c) {Split $W=\ker\pi\oplus W_0$; one coset of $W_0$ captures many graph points};
\draw[arr] (a)--(b);\draw[arr] (b)--(c);
\end{tikzpicture}
\end{center}
\[\pi|_{W_0}\text{ injective}\quad\Longrightarrow\quad
\xi_h=Lh+\xi_0\text{ on a large subset}.\]
\source{\Green{Proposition 1.14 and Proposition 2.6}. This qualitative shortcut is a model, not the quantitative proof above.}
}{
Green's notes offer a simpler finite-field route. The input finite-field
Freiman theorem says: for fixed $p$ and $K$, a finite nonempty
$A\subseteq\F_p^m$ with $|A+A|\leq K|A|$ is contained in a subspace
of size at most $C_{p,K}|A|$. We apply it to $\Gamma'$ in $\F_5^{2n}$.
The projection of the resulting subspace $W$ contains a set of size
$c_\eta|G|$, so $|\ker(\pi|_W)|=|W|/|\pi W|\leq C_\eta$.

Choose a vector-space complement $W_0$ to this kernel. Projection is
injective on each coset of $W_0$; there are only $C_\eta$ cosets inside
$W$. One therefore contains many points of $\Gamma'$. Its inverse
projection is affine, giving $\xi_h=Lh+\xi_0$ on that subset. Extend
$L$ from its subspace to all of $G$ using a basis if needed.

This clarifies the geometry, but these qualitative constants do not
establish the polynomial correlation bound needed in the main theorem.
For those bounds use the graph-refinement/Bogolyubov route, even in
finite fields. Also, a small-radius Bohr set in $\F_p^n$ is an intersection
of kernels, so the corresponding local map is an ordinary linear map.
}

\Backup{Bogolyubov in four Fourier lines}{}{
\small Let $|A|=\delta|G|$, and use normalised convolution.
\[F:=\mathbf1_A*\mathbf1_A*\mathbf1_{-A}*\mathbf1_{-A},
\qquad\supp F\subseteq2A-2A.\]
\[F(x)=\sum_{\xi\in\Gh}|\widehat{\mathbf1_A}(\xi)|^4 e(\xi\cdot x).\]
\[S:=\{\xi:|\widehat{\mathbf1_A}(\xi)|\geq\delta^{3/2}/\sqrt2\},
\qquad |S|\leq2\delta^{-2}.\]
\[x\in\Bohr(S,1/4)\quad\Longrightarrow\quad
F(x)\geq\delta^4-\frac{\delta^3}{2}\sum_\xi|\widehat{\mathbf1_A}(\xi)|^2
=\frac{\delta^4}{2}>0.\]
\tagline{Positive convolution $\Longrightarrow$ an additive representation.}
\source{\GT{Lemma 6.3}.}
}{
This is the proof of the global Bogolyubov statement used in Lemma 2.5.
The transform of the fourfold convolution is the nonnegative fourth power
of the transform magnitude. The zero frequency contributes $\delta^4$.
For every other large frequency, a point in the Bohr set makes the real
part of its character nonnegative. The remaining frequencies have total
fourth-power mass at most $(\delta^3/2)\delta$, by Parseval. The resulting
strictly positive value implies that $x$ has a representation as a sum
of two elements of $A$ minus two others. Parseval also bounds the number
of large frequencies, giving the rank estimate. The subtle local version
in slide 15 replaces global Parseval counting with a local Bessel estimate;
the preparation appendix spells out that replacement.
}

\Backup{local Fourier decay}{}{
\begin{block}{Green--Tao, Lemma 8.4}
Let $U=\Bohr(S,r)$ be regular, $d=\max(1,|S|)$, and
$\ell:\Bohr(S,2r)\to\T$ satisfy
$\ell(x+y)=\ell(x)+\ell(y)$ for $x,y\in U$.
If $|\E_{x\in U}e(\ell(x))|\geq\tau>0$, then
\[\norm{\ell(t)}_{\T}\leq\frac{2^{12}d}{r\tau^2}\norm{t}_S
\qquad(t\in U).\]
\end{block}
\[L\asymp\frac{\tau r}{d\norm{t}_S},\qquad
\frac\tau2\leq\abs{\E_{|j|\leq L}e(j\ell(t))}
\ll\frac1{L\norm{\ell(t)}_{\T}}.\]
\source{\GT{Lemma 8.4}. Regularity gives the first inequality; a geometric series gives the second.}
}{
Take $L$ to be the largest integer with
$L\norm{t}_S\leq\tau r/(400d)$. If $L=0$, the stated bound is trivial
because the right side already exceeds the possible circle distance.
Otherwise the shifts $jt$, $|j|\leq L$, change the mean by at most
$\tau/2$. Local additivity gives
$\ell(x+jt)=\ell(x)+j\ell(t)$, with every intermediate multiple
remaining in the allowed domain. Factor the shifted average and use
that the original mean has magnitude at most one. The normalised
geometric sum in $j$ must have magnitude at least $\tau/2$.
Its usual upper bound gives the displayed conclusion. If
$\norm{t}_S=0$, arbitrary $L$ are allowed, which forces
$\ell(t)=0$ in $\T$. This handles the kernel of the Bohr seminorm.
}

\Backup{integrating on a cyclic group of any order}{}{
\small If $B:(\Z/N\Z)^2\to\T$ is globally bilinear, then
\[B([n],[m])=\frac{a nm}{N}\pmod1\qquad(a\in\Z).\]
\[\boxed{\phi([n])=\frac{a\,n(n-N)}{2N}\pmod1.}\]
\[\phi(n+N)-\phi(n)=an\in\Z\]
\[\phi(n+m)-\phi(n)-\phi(m)=\frac{a nm}{N}\pmod1.\]
\tagline{The linear correction enforces periodicity, even when $N$ is even.}
\source{\JT{proof of Lemma 2.4(i)}, rewritten from its binomial-coefficient formula.}
}{
This optional calculation explains the algebra behind the integration
black box on cyclic groups. Bilinearity forces $B(1,1)$ to be
$N$-torsion in $\T$, so it equals $a/N$. The displayed primitive is
$an^2/(2N)-an/2$. The linear term does not affect the polarisation,
but makes the primitive $N$-periodic: replacing $n$ by $n+N$ adds
the integer $an$. This works for both odd and even $N$.

For $\Z$, use $\alpha n^2/2$. For a product of groups, sum the
primitives on individual factors and add one cross term for each pair
of distinct factors. Symmetry makes this polarise to the full bilinear
form. The structure theorem for finitely generated abelian groups then
proves Lemma 2.4(i). This derivation is for preparation or questions;
the assignment allows the lemma itself to be black-boxed in the talk.
}
